% !TEX TS-program = LuaLaTeX+se

\documentclass[12pt]{book}

%%%%%%%%%%%%%%% STANDARD PACKAGES %%%%%%%%%%%%%%%

%% for personal reasons, the commands are a mix of French and English names. Feel free to choose one or the other if you make use of them. This file is subject to change later on.

\usepackage{geometry}
\usepackage{fontspec}
\usepackage[english,ecclesiasticlatin.usej,activeacute]{babel}
     \babelprovide[hyphenrules=latin]{ecclesiasticlatin}
%          \usepackage{xspace} %% better to use {} after certain macros (or even within definitions) 
     \usepackage{multicol}
\usepackage{paracol}
\footnotelayout{m} %% this allows merged footnote if text is in parallel columns (mostly translations)
\usepackage{fancyhdr}
\usepackage{needspace} %%pour réserver de l'espace en bas de page ; ça évite des séparations entre les titres et la partition ou le texte qui les suivent.
\usepackage[compact]{titlesec}
\usepackage{emptypage}
\usepackage{xcolor}
\usepackage{perpage}%%%to reset footnotes at each page%%%
%\usepackage{xstring}
\usepackage{enumitem} %% nécessaire pour les textes des psaumes.
\usepackage{etoolbox}
\usepackage{lettrine}
%\usepackage{tocloft}%%to fix toc title. titlesec companion packages could also work.
\usepackage{indentfirst}%% to allow for \section with indented 1st paragraphs in good Franco-Latin style (LaTeX default is to suppress this. Package sets  \@afterindentfalse to (always) true.
\usepackage[savepos]{zref}
\usepackage{hyperref}
\usepackage{refcount}


%%%%%%%%%%%%%%% HYPHENATION AND TYPOGRAPHICAL CONVENTIONS %%%%%%%%%%%%%%

%% page settings %%%%
\sloppy %% to avoid overfull hboxes

%\setlength{\textwidth}{5.75in}

\geometry{letterpaper}
\geometry{bindingoffset=0.5in,
inner=1in,
outer=1.35in,
top=1.25in,
bottom=1.25in,
headsep=0.75in
} %%

%% General scale of all graphical elements. Also borrowed from MB
%% Values different from 1 are largely untested.
%% Used in those commands (e.g. everything FontSpec) that use a scale parameter.
\newcommand{\customscale}{1}

\AtBeginDocument{\setlength{\parindent}{1em}} %% should set 1em to EB Garamond not Computer Modern.

%%%%%%%%%%%%%%% GREGORIO CONFIG %%%%%%%%%%%%%%%

\usepackage[autocompile]{gregoriotex}

%% sets * and † to font called via fontspec
\def\GreStar{*}
\def\GreDagger{†}

%%prints asterisk instead of star normally inserted by <v>\greheightstar</v>
 \gredefsymbol{greheightstar}{EB Garamond}{asterisk}

%% should prevent flat turning custos into a clef, as Solesmes  uses an altered custos in only one chant per Élie Roux.
\gresetcustosalteration{invisible}

%produces a score with a smaller initial (default is 40 pt) and annotations like in the Solesmes books; the 1st is optional and can be omitted as needed.
 \NewDocumentCommand{\gscore}{O{}mm}{%
 \greannotation{#1}%
\greannotation{#2}%
\grechangestyle{initial}{\fontsize{28}{28}\selectfont}%
 \gregorioscore{partitions/#3}}
 
 
  \NewDocumentCommand{\lectionscore}{O{}m}{%
   \grecommentary{#1}%
  \grechangestyle{initial}{\fontsize{36}{36}\selectfont}%
  \gregorioscore{partitions/#2}}
 
 %% default should be 40, if not using, annotations need to be turned off in gabc. For the first antiphon of the office, Gospel canticle etc. (initial is both to refer to first and the large illustrated initial, used sometimes with an annotation, cf. Christus factus est of Tenebrae in the ritus monasticus.
  \NewDocumentCommand{\initialscore}{O{}m}{%
      \grechangedim{initialraise}{0.5cm}{scalable}%
  \greannotation{#1}%
   \grechangestyle{initial}{\fontspec{EB Garamond Initials}\fontsize{45}{45}\selectfont}%
 \gregorioscore{partitions/#2}
   \grechangedim{initialraise}{0cm}{scalable}}

%% score with no initial, e.g psalms, common tones, etc.
\newcommand{\smallscore}[2][y]{
  \gresetinitiallines{0}
  \gregorioscore{partitions/#2} %% à remplacer avec NewDocumentCommand ; les arguments ne marchent pas correctement … que fait le [y] ?
  \gresetinitiallines{1}
}

\NewDocumentCommand{\zerobaroffsettextleft}{}%
{%
\grechangedim{maxbaroffsettextleft}%
{0cm}%
{scalable}%
}

\NewDocumentCommand{\zerobaroffsettextright}{}%
{%
\grechangedim{maxbaroffsettextright}%
{0cm}%
{scalable}%
}

%%%the default for left bar offset is 0.3cm. Right is 0.15cm. This needs to be applied after a score where the bar offset is modified.

\NewDocumentCommand{\resetbaroffsettextleft}{}%
{%
\grechangedim{maxbaroffsettextleft}%
{0.3cm}%
{scalable}%
}

\NewDocumentCommand{\resetbaroffsettextright}{}%
{%
\grechangedim{maxbaroffsettextright}%
{0.15cm}%
{scalable}%
}

 \NewDocumentCommand{\MagAnnotation}{}{%
 \grechangedim{annotationseparation}{-0.01cm}{fixed}%
 } %% pour les antiennes où le chiffre convient mieux aux miniscules (parfois 1, 2, 4, 7…) ; il n'est pas nécessaire pour 6 ou 8.
%\gresetheadercapture{annotation}{greannotation}{string}

%% A-bove L-ines T-ext shall be italicized and in a smaller font
\grechangestyle{abovelinestext}{\itshape\fontsize{8}{10}\selectfont}
%% for psalm tones etc. Roman/upright text is what the Liber Usualis uses. 
\NewDocumentCommand{\altnormal}{}{\grechangestyle{abovelinestext}{\normalfont\fontsize{8}{10}\selectfont}}
%%  resets ALT font
\NewDocumentCommand{\altitshape}{}{\grechangestyle{abovelinestext}{\itshape\fontsize{8}{10}\selectfont}} 

\NewDocumentCommand{\altlarge}{}{\grechangestyle{abovelinestext}{\large}} %% for I reading of Lamentations

\NewDocumentCommand{\altupright}{}{\grechangestyle{abovelinestext}{\normalfont}} %%


%% fine-tuning of space beween the staff and the text above lines (used for Magnificats; needs to be rethought (again, should 2 and 3 be raised or not?) and worked into \smallscore
\newcommand{\altraise}{0.1cm} %% default is -0.1cm %% needs to be fine-tuned for \rubrique command with EB Garamond font. -0.2cm is MB value
\grechangedim{abovelinestextraise}{\altraise}{scalable}

\newcommand{\altheight}{0.5cm} %% default is 0.3cm and value must be bigger than text height; this should fix the problem. but needs further investigation.
\grechangedim{abovelinestextheight}{\altheight}{scalable} %% this is a very finicky command.

\newcommand{\comraise}{0.4cm} %% default is 0,2m %%
\grechangedim{commentaryraise}{\comraise}{scalable}

%% allows printing of glyphs from Gregorio score font (available glyphs listed in documentation) as text.
\makeatletter
\def\gretextglyph#1{{\gre@font@music\csname GreCP#1\endcsname}}
\makeatother

%%% Font %%%
\setmainfont{EB Garamond}[UprightFont=EB Garamond Regular,
ItalicFont= EB Garamond Italic,
BoldFont= EB Garamond Bold,
Ligatures=Rare,
Numbers=Proportional,
Numbers=OldStyle] %% all \kern numbers are based on this and can be removed if this font is abandoned.
\newfontfamily\symbolfont{liturgy}
  %% text must be written {\symbolfont{+}} (you can use V and R as well) but may conflict as + is † in gabc. Cross should be \small or even \footnotesize


%    \grechangestyle{initial}{\fontspec{EB Garamond Initials}\fontsize{#1}{#1}\selectfont}

%%% default initial size is 32 points
%\newcommand{\defaultinitialsize}{28}
%\initialsize{\defaultinitialsize}

%%%% Font %%%
%\setmainfont{EB Garamond}[UprightFont=EB Garamond Regular,
%ItalicFont= EB Garamond Italic,
%BoldFont= EB Garamond Bold,
%Ligatures=Rare,
%StylisticSet=6,
%Numbers=Proportional,
%Numbers=OldStyle] %% all \kern numbers are based on this and can be removed if this font is abandoned. %%StylisticSet=6 is, in Pardo EBGaramond, the long-tailed Q.
%\newfontfamily\symbolfont{liturgy} 
  %% text must be written {\symbolfont{+}} (you can use V and R as well) but may conflict as + is † in gabc. Cross should be \small or even \footnotesize
  
  %%%Roman numerals for the year %% https://tex.stackexchange.com/questions/185548/the-year-in-roman-and-the-month-in-text

\makeatletter
\newcommand{\YEAR}{\@Roman{\the\year}}
\makeatother


%%% TYPOGRAPHICAL AND SYMBOL COMMANDS  %%%%

%% \P is pilcrow and is defined in LaTeX as is.
%% Add {\textcolor{gregoriocolor} as first part of command's definition if changing to red.


%%for proper spacing of all-caps letters to prevent clashes, e.g. serif strokes touching.

\NewDocumentCommand\capspace{m}{{\addfontfeature{LetterSpace=5.0}{#1}}}
\NewDocumentCommand\scspace{m}{\textsc{{{\addfontfeature{LetterSpace=5.0}{#1}}}}}

\NewDocumentCommand\feasttitle{m}{{\centering{\capspace{\huge{#1}}}\par}}


%% macro to print Alleluia for versicles.
\NewDocumentCommand{\tpalleluia}{}{(\textit{T.P.} Allelúia.)}

\newcommand{\specialcharhsep}{3mm} % space after invoking R/ or V/ or A/ outside rubrics
%\newcommand{\vv}{\textcolor{gregoriocolor}{\fontspec[Scale=\customscale]{Charis SIL}℣.\nolinebreak[4]\hspace{\specialcharhsep}\nolinebreak[4]}} %% format from MB

\newcommand{\vv}{{\normalfont ℣.\nolinebreak[4]\hspace{\specialcharhsep}\nolinebreak[4]}}
\newcommand{\rr}{{\normalfont ℟.\nolinebreak[4]\hspace{\specialcharhsep}\nolinebreak[4]}}

%% typesets a cross pattée.
\NewDocumentCommand{\cc}{}{%
    % Grouping to keep font changes local
    {%
        % Ensure we're not in italics (since liturgy.ttf doesn't have italic)
        \normalfont%
        % Select the font liturgy.ttf
        \symbolfont%
        % Set the size
        \footnotesize%
        % In liturgy.ttf, the plus (+) is a cross pattée
        +%
    }%    
} %%  can replace <+> with <U+2720> (see above) if a suitable font is found (or EB Garamond font is fixed); command will be useful in lieu of typing the Unicode.
%% use <v> tag to insert into a gabc score (usually for the faithful or for blessings, e.g. the font.

%% Same special characters, for in-score use (<sp>V/ R/ A/ +</sp>)

\gresetspecial{V/}{{\fontspec[Scale=\customscale]{EB Garamond}℣.~}}
\gresetspecial{R/}{{\fontspec[Scale=\customscale]{EB Garamond}℟.~}}
\gresetspecial{+}{{\fontspec[Scale=\customscale]{EB Garamond}†~}}
\gresetspecial{*}{\gresixstar}
\gresetspecial{cross}{\cc}

%% Same special characters, for use in rubrics (no space)

\newcommand{\vvrub}{{\normalfont ℣.~}}
\newcommand{\rrrub}{{\normalfont ℟.~}}

\NewDocumentCommand{\antiphona}{mm}
{\begin{otherlanguage*}{#1}%
\noindent%
 Ant. #2.%
 \end{otherlanguage*}%
 }
 


%%rubrics: black italics, smaller than body of psalms etc

\NewDocumentCommand{\rubrique}{m}{%
    {%
        \fontsize{10}{12}%
        \selectfont%
        \textit{%
        %
        #1%
    }%
    }}
    
%% macro to print normal text inside of rubric (name of a chant or prayer, etc.)

\NewDocumentCommand{\normaltext}{m}{%
    {%
        \normalfont%
        #1%
    }%
    }
    
%% in case something should be bolded inside of a rubrique
\NewDocumentCommand{\rubriquegras}{m}{%
    {%
        \normalfont%
\textbf{#1%
}%
}}

%% to print in red instead of italicizing
\NewDocumentCommand{\rouge}{m}{%
\textcolor{gregoriocolor}%
{\fontsize{10}{12}%
\selectfont{%
#1%
}%
}}

%% to print in black within rouge group.

%\NewDocumentCommand{\textenoir}{m}{%

\NewDocumentCommand{\textenoir}{m}{%
\textcolor{black}%
{%
\normalfont%
#1%
}%
}

%% rouge but in italic (this is technically not correct).
\NewDocumentCommand{\rougeit}{m}{%
\textcolor{gregoriocolor}%
{\fontsize{10}{12}%
\selectfont%
\textit{#1%
}%
}}

%%for Scriptural references of the chapter.
\NewDocumentCommand{\scripture}{m}{%
{\raggedleft{\textit{#1}%
}%
\par}%
}

%\newcommand{\rubriquegras}[1]{{\fontsize{9}{11}\selectfont\textbf{#1}}}

%%macro to space punctuation with ecclesiasticlatin language from babel.
\NewDocumentCommand{\espaceponctuation}{}{\hspace{0.10em}} %%this value seems more balanced with this font than the definition provided in the ecclesiasticlatin documentation.

\NewDocumentCommand{\myrule}{}{%
    \par%
    {%
        \centering%
        \rule{0.3\textwidth}{0.4pt}%
        \par%
    }% typesets a horizontal rule like on the page with the prayers before and after the office.
}  %%If you're using color at all in your document, you might want to either force \myrule to use black or make it cusotmizable. (from u/Independent-Comb-257)

\NewDocumentCommand{\bigrule}{}{%
    \par%
    {%
        \centering%
        \rule{0.6\textwidth}{0.4pt}%
        \par%
    }% typesets a horizontal rule like on the page with the prayers before and after the office.
}  %%If you're using color at all in your document, you might want to either force \myrule to use black or make it cusotmizable. (from u/Independent-Comb-257)


%%% typesets a cross pattée.

\NewDocumentCommand{\mycross}{}{%
{\symbolfont\small{{+}}}%
{}
} %% can replace <+> with <U+2720> if a suitable font is found (or EB Garamond font is fixed); command will be useful in lieu of typing the Unicode.

%% macro to format psalm text

%\NewDocumentCommand{\pstexte}{ m }{%
%    \smallskip%
%    \noindent%
%    \begin{itemize}[%
%    		label=\null, %
%			leftmargin=0pt, %
%			itemindent=0mm, %
%			labelsep=0pt, %
%			labelwidth=0pt, %
%			rightmargin=0pt, %
%			parsep=0pt, %
%			topsep=0pt, %
%			itemsep=0pt]%
%    \input{psaumes/#1}
%    \end{itemize}}
          
      \setlength{\columnsep}{3pc}
%\setlength{\columnsep}{10pt}
\setlength\columnseprule{0.4pt}


%% original version kept so that I can fix errors found later and keep using common texts, going forward
    \NewDocumentCommand{\textebilingue}{ mm }{%
    \smallskip%
       \begin{paracol}{2}
           \input{psaumes/#1}
           \switchcolumn
    \input{psaumes/#2}
    \end{paracol}}

%2nd version;  no need for input (too many chapter/collect files) + is the equivalent of long, allows \par
%    \NewDocumentCommand{\texteparacol}{ +m+m }{%
%    \smallskip%
%       \begin{paracol}{2}
%           #1
%           \switchcolumn
%    #2
%    \end{paracol}}

%%%3rd version: need to switch languages!
    \NewDocumentCommand{\texteparacol}{ +mm+m }{%
    \smallskip%
       \begin{paracol}{2}
           #1
           \switchcolumn
           \begin{otherlanguage}{#2}
    #3
    \end{otherlanguage}
    \end{paracol}}
    
    
    %%psalm vernacular file has language command
    \NewDocumentCommand{\pstexte}{mm}{%
    \smallskip%
    \noindent%%
   \begin{paracol}{2}
    \begin{enumerate}[%
    		label=\arabic*., %
			leftmargin=3pt, %
			itemindent=3mm, %
			labelsep=3pt, %
			labelwidth=0pt, %
			rightmargin=0pt, %
			parsep=0pt, %
			topsep=0pt, %
			itemsep=0pt,%
			start=2]%
    \input{psaumes/#1.tex}
    \end{enumerate}
    \switchcolumn
       \begin{enumerate}[%
    		label=\arabic*.,%
			leftmargin=3pt, %
			itemindent=3mm, %
			labelsep=3pt, %
			labelwidth=0pt, %
			rightmargin=0pt, %
			parsep=0pt, %
			topsep=0pt, %
			itemsep=0pt]%
    \input{psaumes/#2.tex}
      \end{enumerate}
    \end{paracol}%
    } %% modified from original in order to allow 2 column psalms
    
    %% Command de Matthias Bry modifié, prints psalm incipit score with the text
    \NewDocumentCommand{\psalmus}{mmmm}{%
	\needspace{4\baselineskip}%
	\smalltitle{Psalmus #1.}%
	\smallscore[n]{#2}%
	\pstexte{#3}{#4}%
}

    %% sd=semiduplex. especially for Dominica ad Vesperas in the Psalterium (the only example which comes to mind, in fact!)
 \NewDocumentCommand{\sdpsalmusbilingue}{mmm}{%
	\needspace{4\baselineskip}%
		\smalltitle{Psalmus #1.}%
	\pstexte{#2}{#3}%
	}

    %% Command de Matthias Bry modifié, prints canticle incipit score with the text and scriptural commentary
    \NewDocumentCommand{\canticum}{mmmm}{
	\needspace{4\baselineskip}
	\smalltitle{Canticum #1.}
	 \grecommentary{#2}
	\smallscore[n]{#3}
	\pstexte{#4}
}
	
    %% macro to print any additional text (Capitiulum, oratio, rubrics)
 \NewDocumentCommand{\textes}{ m }{%
    \input{textes/#1}}
    
%     \NewDocumentCommand{\lection}{m}{%
%   \begin{multicols}{2}#1\end{multicols}} %% this doesn't work as expected, needs to be fixed
    
    
    
    %% SC work for page headers and for secondary headings but not so much for things like "Dominica…" and certainly not the feast name%%
    
    %%%% Headers%%%%
    \pagestyle{fancy}
\fancyhead{}
\fancyfoot{}
\renewcommand{\headrulewidth}{0pt}


 \setlength{\headheight}{13.59999pt} %% recommended by LaTeX

%\fancyhead[RO]{\small\thepage}
%\fancyhead[LE]{\small\thepage}

\fancyhead[RO]{\small\rightmark\hspace{0.5cm}\thepage}
\fancyhead[LE]{\small\thepage\hspace{0.5cm}\leftmark}

\newcommand{\setheaders}[2]{
	\renewcommand{\rightmark}{{\sc#2}}
	\renewcommand{\leftmark}{{\sc#1}}
}
\setheaders{}{}


%% TITLE Commands %%%% %% TITLE Commands %%%% currently dead but Matthias says this isn't a good idea.
 \titleformat{\chapter}[display]{\Huge\filcenter\sc\addfontfeature{LetterSpace=5.0}}{}{0pt}{} %%
\titleformat{\section}[display]{\LARGE\filcenter\sc\addfontfeature{LetterSpace=5.0}}{}{}{} %%
\titleformat{\subsection}[display]{\Large\filcenter\sc\addfontfeature{LetterSpace=5.0}}{}{}{} %%originally \large
\setcounter{secnumdepth}{0}
\titlespacing*{\chapter}{0pt}{-25pt}{20pt} %%this should reduce spacing before chapter title while putting it flush with the top of the text area
%% see https://tex.stackexchange.com/questions/63390/how-to-decrease-spacing-before-chapter-title
%%\titlespacing*{\chapter}{0pt}{-50pt}{40pt}  %%this puts chapter title in HEADER which we do not want
\titlespacing{\section}{0pt}{*0}{*0}

%%need section titleformat more appropriate for weekdays of psalter and minor feasts 
% See https://tex.stackexchange.com/questions/623797/multiple-section-styles-in-same-document.

\addto\captionsecclesiasticlatin{\renewcommand\contentsname{\centerline{INDEX GENERALIS.}}}

% \NewDocumentCommand{\mysection}{m}{\section{#1}\setheaders{{\scshape\addfontfeature{LetterSpace=5.0}#1}}{{\scshape\addfontfeature{LetterSpace=5.0} #1}}}
 
\NewDocumentCommand{\header}{m}{\setheaders{{\scshape\addfontfeature{LetterSpace=5.0}#1}}{{\scshape\addfontfeature{LetterSpace=5.0} #1}}}

%% all of these are a mess and should probably be ignored, but feel free to use them by referring to Matthias B's originals in the Nocturnale Romanum repository.

\newcommand{\officiumtitulum}[1]{
%  \newpage
%  \thispagestyle{empty}
  \setheaders{{\scshape\addfontfeature{LetterSpace=3.0}#1}}{{\scshape\addfontfeature{LetterSpace=3.0} #1}}
  \begin{center}
  {\scshape\large #1}\par
  \end{center}}
  
%   \setheaders{{{\scshape\addfontfeature{LetterSpace=3.0}}}{{{\scshape\addfontfeature{LetterSpace=3.0} {#1}}}

%%formatting date of feasts where this listed at top of page
\NewDocumentCommand{\litdate}{m}{%
\smalltitle{Die #1.}%
}

\NewDocumentCommand{\litdateen}{m}{%
\smalltitle{#1.}%
}

%%formatting for Sunday feasts : inter, post… name is for consistency
\NewDocumentCommand{\dominicadate}{m}{%
\smalltitle{Dominica #1.}%
}

\NewDocumentCommand{\sundaydate}{m}{%
\smalltitle{Sunday #1.}%
}

%%will use litdaten for English
%%formatting for feasts that fall on a weekday fixed to some other thing, like Corpus Christi, Sacred Heart, BVM in Passiontide, and the new feasts found in the appendix (Feria Roman numeral in Latin: inter, post… name is for consistency
\NewDocumentCommand{\feriadate}{m}{%
\smalltitle{Feria #1.}%
}

 %%looks like xparse formatting changed so o just has No Default Value whereas O{} is that OR you can insert something in braces!
%%answer thanks to David Carlisle
\NewDocumentCommand{\duplexclassis}{mo}{%
\needspace{5\baselineskip}%
\vspace{\baselineskip}%
 {\centering\textit{\large Duplex #1 classis\IfValueT{#2}{ #2}.}\par}%
}

\NewDocumentCommand{\doubleclass}{mo}{%
\needspace{5\baselineskip}%
\vspace{\baselineskip}%
 {\centering\textit{\large Double of the #1 class\IfValueT{#2}{ #2}.}\par}%
}


\NewDocumentCommand{\duplexmajus}{}{%
\needspace{5\baselineskip}%
\vspace{\baselineskip}%
 {\centering{\textit{\large Duplex majus.}}\par}%
}

\NewDocumentCommand{\doublemajor}{}{%
\needspace{5\baselineskip}%
\vspace{\baselineskip}%
 {\centering{\textit{\large Double major.}}\par}%
}

%%is Duplex even specified.as it's the default?
\NewDocumentCommand{\duplex}{}{%
\needspace{5\baselineskip}%
\vspace{\baselineskip}%
 {\centering{\textit{\large Duplex.}}\par}%
}

\NewDocumentCommand{\double}{}{%
\needspace{5\baselineskip}%
\vspace{\baselineskip}%
 {\centering{\textit{\large Double.}}\par}%
}

\NewDocumentCommand{\semiduplex}{}{%
\needspace{5\baselineskip}%
\vspace{\baselineskip}%
 {\centering{\textit{\large semiuplex.}}\par}%
}

\NewDocumentCommand{\semidouble}{}{%
\needspace{5\baselineskip}%
\vspace{\baselineskip}%
 {\centering{\textit{\large Semidouble.}}\par}%
}

\NewDocumentCommand{\simplex}{}{%
\needspace{5\baselineskip}%
\vspace{\baselineskip}%
 {\centering{\textit{\large Simplex.}}\par}%
}

\NewDocumentCommand{\simple}{}{%
\needspace{5\baselineskip}%
\vspace{\baselineskip}%
 {\centering{\textit{\large Simple.}}\par}%
}

\NewDocumentCommand{\rank}{m}{
\needspace{5\baselineskip}  %% very small if there are neumes above the staff, including flats in mode 2, e.g. O Doctor optime
\vspace{\baselineskip}
 {\centering{\textit{\large #1}}\par}
}



\NewDocumentCommand{\bigtitle}{m}{
\needspace{5\baselineskip}  %% very small if there are neumes above the staff, including flats in mode 2, e.g. O Doctor optime
\vspace{\baselineskip}
 {\centering{\large#1}\par}
}

\NewDocumentCommand{\smalltitle}{m}{
\needspace{5\baselineskip}  %% very small if there are neumes above the staff, including flats in mode 2, e.g. O Doctor optime
\vspace{\baselineskip}
 {\centering #1\par}
}

\NewDocumentCommand{\chapterlateng}{}{%
\smalltitle{Chapter.}%
}


%\newcommand{\subtitulum}[1]{
% \subsection{
%  \begin{center}
%  {\large#1}\par
%  \end{center}}}
  
  %% something has broken here
  %% see https://tex.stackexchange.com/questions/545841/paragraph-ended-before-ttlformatsi-was-complete

%% originally used apparently for typesetting portions of the Common of the BVM
\newcommand{\espacetitre}{\vspace{-0.5cm}}
\newcommand{\espaceps}{\vspace{-3mm}} %% ps =psaume
\newcommand{\espacecap}{\vspace{-3mm}} %% cap is capitulum
\newcommand{\espace}{\vspace{2mm}}

\NewDocumentCommand{\secreto}{}{Pater noster. \rubrique{secreto.}}

\NewDocumentCommand{\pateravecredo}{}{\large{Pater noster, Ave María, \textit{et} Credo.}}

\NewDocumentCommand{\pateravedeus}{}{%
Pater noster et Ave María. ℣. Deus in adjutórium.%
}

\NewDocumentCommand{\prayers}{}{%
Our Father and Hail Mary, \rubrique{silently.}%
}

\NewDocumentCommand{\orationes}{}{%
Pater noster et Ave María \rubrique{secreto.}%
}

%%for litanies
\makeatletter  
\newcounter{score}
\newcounter{tabstop}[score]
\newcommand{\grealign}{%
	\@bsphack%
	\ifgre@boxing\else%
		\kern\gre@dimen@begindifference%
		\stepcounter{tabstop}%
		\expandafter\zsavepos{stop-\thescore-\thetabstop}%
		\kern-\gre@dimen@begindifference%
	\fi%
	\@esphack%
}

\newcommand{\setstops}{%
  \gdef\nstabbing@stops{%
    \hspace*{-\oddsidemargin}\hspace{-1in}%
    \hspace*{\zposx{stop-\thescore-1} sp}\=%
  }%
  \count@=\@ne
  \loop\ifnum\count@<\value{tabstop}%
    \begingroup\edef\x{\endgroup
      \noexpand\g@addto@macro\noexpand\nstabbing@stops{%
        \noexpand\hspace{-\noexpand\zposx{stop-\thescore-\the\count@} sp}%
        \noexpand\hspace{\noexpand\zposx{stop-\thescore-\the\numexpr\count@+1} sp}\noexpand\=%
      }%
    }\x
    \advance\count@\@ne
  \repeat
  \nstabbing@stops\kill
}
\makeatother

\newenvironment{nstabbing}
  {\setlength{\topsep}{0pt}%
   \setlength{\partopsep}{0pt}%
   \tabbing%
   \setstops}
  {\endtabbing\stepcounter{score}}

\makeatletter
\newcounter{blankpages} %% should remove number on blank page at end of preface. Cf https://tex.stackexchange.com/questions/250761/how-to-not-count-blank-pages-in-frontmatter-of-two-sided-document
\def\cleardoublepage{%
    \clearpage
    \if@twoside
    \ifodd\c@page
    \else
    \hbox{}
    \thispagestyle{empty}%
    \newpage\stepcounter{blankpages}%
    \if@twocolumn\hbox{}\newpage\fi
    \fi
    \fi
}
\newcommand{\@romannoblank}[1]{%
    \@roman{\numexpr#1-\value{blankpages}\relax}%
}
\makeatother

\providecommand\phantomsection{} %% should make hyperref work properly %%only works with section etc., not with label etc. (that requires \phantomsection
%

\NewDocumentEnvironment{enpars}{}
  {\begin{otherlanguage*}{english}}
  {\par\end{otherlanguage*}}
  
  \NewDocumentEnvironment{frpars}{}
  {\begin{otherlanguage*}{french}}
  {\par\end{otherlanguage*}}
%\renewcommand{\contentsname}{\hfill\bfseries\LARGE Index Generalis.\hfill\null}  %% modifies TOC title
%\renewcommand{\cfttoctitlefont}{\hfil\LARGE}

%%%SUBFILES%%%
\usepackage{xr} %%to allow cross-references between documents%%%
  \usepackage{subfiles}
  
  %% When we start a new subfile (new chapter or section) , 
%% we start on a new page (with blank filling on the previous page) and create a corresponding label.

\newcommand{\customsubfile}[1]{\newpage\label{#1}\thispagestyle{empty}\subfile{#1}}


\begin{document}
 
 \twosided[pc] %% to allow for the left-right switching of columns needed for bilingual liturgical typesetting (like in missals); it's unclear why this goes BEFORE everything else and not when activating the environment. %% It's OK for paper booklets to not switch, but it's probably best practice to switch.
 
 \thispagestyle{empty}
 
 {\centering{\huge{SUNDAYS OF LENT AT VESPERS.}}\par}\header{sundays of lent at vespers.}
 
\smalltitle{Before the Office.}

\rubrique{These prayers are said kneeling. End when the priest stands.}

\textebilingue{orationes_ante_do}{orationes_ante_do_eng}

\smalltitle{\prayers}

\rubrique{The celebrant intones the following verse. All sing the response and doxology.}

\begin{otherlanguage}{english}
\rubrique{Make the sign of the cross at the beginning and at the Magnificat, and bow the head when the Trinity is mentioned, when mention is made of the ``Holy Name'' in some way, and at ``Orémus.''}\end{otherlanguage}

\gscore[]{℣.}{Deus_In_Adjutorium_laus_tibi}

\noindent
\begin{otherlanguage}{english}
O God, \cc{} come to my assistance. \rr O Lord, make haste to help me. Glory be to the Father, and to the Son,~* and to the Holy Ghost. As it was in the beginning, is now,~* and ever shall be, world without end. Amen. Praise be unto Thee, O Lord, King of Eternal Glory.\end{otherlanguage}

\smalltitle{First Antiphon. 7. c2.}
\initialscore{1_Dixit_Dominus}

\footnotetext[1]{The words are not repeated, which is always done when the antiphon begins with the first words of the psalm, which continues therefore after the word with which the antiphon ends, whether it was sung in full or simply begun. Regardless, if the words are the same, the antiphon is continued with the psalm.}

\rubrique{\begin{otherlanguage}{english}
Sit at the word ``Sede.'' The cantor sings odd verses; all sing even verses, and this for the hymn and Magnificat as well.
\end{otherlanguage}}

\sdpsalmusbilingue{109}{109_7c2}{109_en}

\smallscore{an--dixit_dominus_domino--solesmes_1961}

\begin{otherlanguage}{english} \noindent Ant. The Lord said to my Lord: Sit thou at my right hand. \end{otherlanguage}

\rubrique{\begin{otherlanguage}{english}
Clerics in descending order, or a cantor, intone the subsequent antiphons. The cantor then continues the psalm.
\end{otherlanguage}}

\gscore[2. Ant.]{3. b}{2_Magna_Opera}

\psalmus{110}{110_3b}{110_3b}{110_en}

\smallscore{an--magna_opera_domini--solesmes_1961}

\begin{otherlanguage}{english} \noindent Ant. Great are the works of the Lord, sought out according to all his wills. \end{otherlanguage}

\gscore[3. Ant.]{4. g}{3_Qui_timet_Dominum}

\psalmus{111}{111_4g}{111_4g}{111_en}

\smallscore{an--qui_timet_dominum--solesmes_1961}

\begin{otherlanguage}{english} \noindent Ant. He who feareth the Lord shall delight exceedingly in his commandments. \end{otherlanguage}

\gscore[4. Ant.]{7. c}{4_Sit_nomen_Domini}

\psalmus{112}{112_7c}{112_7c}{112_en}

\smallscore{an--sit_nomen_domini..._in--solesmes_1961}

\begin{otherlanguage}{english} \noindent Ant. Blessed be the name of the the Lord forever. \end{otherlanguage}

\gscore[5. Ant.]{T. pereg.}{5_Deus_autem_noster}

\psalmus{113}{113_per}{113_per}{113_en}

\smallscore{an--deus_autem_noster--solesmes_1961}

\begin{otherlanguage}{english} \noindent Ant. But our God is in heaven, he hath done all things whatsoever he would. \end{otherlanguage}


\rubrique{\begin{otherlanguage}{english}
Having sung the last antiphon in full, stand when the celebrant does.
\end{otherlanguage}}

%\newpage

\bigtitle{Chapter.}

\smalltitle{I Sunday of Lent.}

\scripture{2 Cor. 6, 1 – 2.}
\texteparacol{\lettrine{F}{r}atres: Hortámur vos, ne in vácuum grátiam Dei recipiátis.~† Ait enim: Témpore accépto exaudívi te,* et in die salútis adjúvi te.

\rr Deo grátias.}{\noindent Brothers, we exhort you, that you receive not the grace of God in vain. For he saith: In an accepted time have I heard thee; and in the day of salvation have I helped thee. Thanks be to God.}

\smalltitle{II Sunday of Lent.}

\scripture{1 Thess. 4, 1.}

\texteparacol{\lettrine{F}{r}atres: Rogámus vos, et obsecrámus in Dómino Jesu:~† ut, quemádmodum accepístis a nobis, quómodo vos opórteat ambuláre, et placére Deo,* sic et ambulétis, ut abundétis magis.

\rr Deo grátias.}{\noindent For the rest therefore, brethren, we pray and beseech you in the Lord Jesus, that as you have received from us, how you ought to walk, and to please God, so also you would walk, that you may abound the more. Thanks be to God.}

\smalltitle{III Sunday of Lent.}

\scripture{Eph. 5, 1 – 2.}
\texteparacol{\lettrine{F}{r}atres: Estóte imitatóres Dei, sicut fílii caríssimi:~† et ambuláte in dilectióne, sicut et Christus diléxit nos, et trádidit semetípsum pro nobis~* oblatiónem et hóstiam Deo in odórem suavitátis.

\rr Deo grátias.}{\noindent Brothers: Be ye, therefore, followers of God, as most dear children; And walk in love, as Christ also hath loved us, and hath delivered himself for us, an oblation and a sacrifice to God for an odour of sweetness. Thanks be to God.}

\smalltitle{IV Sunday of Lent.}

\scripture{Gal. 4, 22 – 24.}
\texteparacol{\lettrine{F}{r}atres: Scriptum est quóniam Abraham duos fílios hábuit: unum de ancílla, et unum de líbera:~† sed qui de ancílla, secúndum carnem natus est: qui autem de líbera, per repromissiónem:~* quæ sunt per allegoríam dicta.

\rr Deo grátias.}{\noindent Brothers: For it is written that Abraham had two sons: the one by a bondwoman, and the other by a free woman. But he who was of the bondwoman, was born according to the flesh: but he of the free woman, was by promise, which things are said by an allegory. Thanks be to God.}

\smalltitle{Passion Sunday.}

\scripture{Heb. 9, 11 – 12.}

\texteparacol{\lettrine{F}{r}atres: Christus assístens Póntifex futurórum bonórum, per ámplius et perféctius tabernáculum non manu factum, id est, non huius creatiónis:~† neque per sánguinem hircórum aut vitulórum, sed per próprium sánguinem introívit semel in Sancta,~* ætérna redemptióne invénta.

\rr Deo grátias.}{\noindent Brothers: But Christ, being come an high priest of the good things to come, by a greater and more perfect tabernacle not made with hand, that is, not of this creation: Neither by the blood of goats, or of calves, but by his own blood, entered once into the holies, having obtained eternal redemption. Thanks be to God.}

\smalltitle{Palm Sunday.}

\scripture{Phil. 2, 5 – 7.}
\texteparacol{\lettrine{F}{r}atres: Hoc enim sentíte in vobis, quod et in Christo Jesu: qui, cum in forma Dei esset, non rapínam arbitrátus est esse se æquálem Deo:~† sed semetípsum exinanívit, formam servi accípiens, in similitúdinem hóminum factus,~* et hábitu invéntus ut homo.

\rr Deo grátias.}{\noindent Brothers: For let this mind be in you, which was also in Christ Jesus: Who being in the form of God, thought it not robbery to be equal with God: But emptied himself, taking the form of a servant, being made in the likeness of men, and in habit found as a man. Thanks be to God.}

%% The usual tone and the 3rd tone, used ad libitum, are omitted as we do not use these tones, but they are included in the tex file and the gabc files are available, for the sake of completeness.

\smalltitle{Hymn.}

\rubrique{This hymn is sung at Vespers of Lent, from Saturday before the I Sunday of Lent until the Saturday before Passion Sunday exclusive.}

%\gscore[]{8.}{hy--Lucis_Creator_optime}
%
%\rubrique{\begin{otherlanguage}{english}
%If, however, a commemoration is made of the Blessed Virgin Mary:
%\end{otherlanguage}}
%
%\smallscore{hy--Lucis_Creator_BMV}
%
%\smalltitle{2. Another tone ad lib.}

\gscore[]{2.}{hy_audi_benigne_conditor_solesmes_1961}

\begin{multicols}{2}
\begin{otherlanguage}{english}
O kind Creator, bow thine ear\\
To mark the cry, to know the tear\\
Before thy throne of mercy spent\\
In this thy holy fast of Lent.

Our hearts are open, Lord, to thee:\\
Thou knowest our infirmity;\\
Pour out on all who seek thy face\\
Abundance of thy pardoning grace.

Our sins are many, this we know;\\
Spare us, good Lord, thy mercy show;\\
And for the honour of thy name\\
Our fainting souls to life reclaim.

Give us the self-control that springs\\
From discipline of outward things,\\
That fasting inward secretly\\
The soul may purely dwell with thee.

We pray thee, Holy Trinity,\\
One God, unchanging Unity,\\
That we from this our abstinence\\
May reap the fruits of penitence. Amen.
\end{otherlanguage}
\end{multicols}

\rubrique{\begin{otherlanguage}{english}
The cantor sings the verse to which all reply.
\end{otherlanguage}}

\smallscore{versiculus_quadragesimae}

\begin{otherlanguage}{english} \noindent \vv  God hath given His Angels charge over thee.

\noindent \rr To keep thee in all thy ways.
 \end{otherlanguage}
 
 \smalltitle{Hymn.}
 
 \rubrique{This hymn is sung at Vespers of Passiontide, from Saturday before Passion Sunday until Wednesday of Holy Week inclusive.}

\rubrique{All kneel for the entire strophe beginning \normaltext{O crux ave Spes única.}}
 
 \gscore[]{1.}{hy_vexilla_regis_prodeunt_solesmes}
 
 \smallscore{versiculus_passionis}

\begin{otherlanguage}{english} \noindent \vv Deliver me, O Lord, from the evil man.

\noindent \rr Rescue me from the unjust man.
 \end{otherlanguage}
 
 \rubrique{\begin{otherlanguage}{english}
The celebrant intones the \normaltext{Magnificat} antiphon; this is always proper to the day. The cantor continues the \normaltext{Ma\-gnificat.}\end{otherlanguage}}

\bigtitle{Canticle of the Blessed Virgin Mary.}

\smalltitle{I Sunday of Lent.}\label{1_sunday}

\gscore[]{Ant. 8. G*}{an_ecce_nunc_tempus_solesmes_1961}
\begin{otherlanguage}{english} \noindent Ant. Behold, now is the accepted time behold, now is the day of salvation; in these days therefore let us approve ourselves as the ministers of God, in much patience, in fastings, in watchings, and in love unfeigned. \end{otherlanguage}

\gresetinitiallines{0}

\gabcsnippet{(c4) 1. Ma(g)gní(hg)fi(gj)cat(j.) *(:) á(j)ni(j)ma(j) <i>me</i>(i)<i>a</i>(j) <b>Dó</b>(h)mi(gr)num.(gh..) (::)}

\gabcsnippet{(c4) 2. Et(g) ex(h)sul(j)tá(j)vit(j) spí(j)ri(j)tus(j) <b>me</b>(k//jr)us(j.) *(:) in(j) De(j)o(j) sa(j)lu(j)<i>tá</i>(i)<i>ri</i>(j) <b>me</b>(h//gr)o.(gh..) (::)}

\gresetinitiallines{1}
   \smallskip%
    \noindent%
    \begin{itemize}[%
    		label=\null, %
			leftmargin=0pt, %
			itemindent=0mm, %
			labelsep=0pt, %
			labelwidth=0pt, %
			rightmargin=0pt, %
			parsep=0pt, %
			topsep=0pt, %
			itemsep=0pt]%

\item 3. Quia respéxit humilitátem ancíllæ \textbf{su}æ:~* ecce enim ex hoc beátam me dicent omnes gene\textit{rati}\textbf{ó}nes.

\item 4. Quia fecit mihi magna qui \textbf{pot}ens est:~* et sanctum \textit{nomen} \textbf{e}jus.

\item 5. Et misericórdia ejus a progénie in pro\textbf{gé}nies~* timén\textit{tibus} \textbf{e}um.

\item 6. Fecit poténtiam in bráchio \textbf{su}o:~* dispérsit supérbos mente \textit{cordis} \textbf{su}i.

\item 7. Depósuit poténtes de \textbf{se}de,~* et exal\textit{távit} \textbf{hú}miles.

\item 8. Esuriéntes implévit \textbf{bo}nis:~* et dívites dimí\kern 0.04em\textit{sit} \textit{in}\textbf{á}nes.

\item 9. Suscépit Israel púerum \textbf{su}um,~* recordátus misericór\textit{diæ} \textbf{su}æ.

\item 10. Sicut locútus est ad patres \textbf{no}stros,~* Abraham et sémini e\kern 0.02em\textit{jus} \textit{in} \textbf{sǽ}cula.

\item 11. Glória Patri, et \textbf{Fí}lio,~* et Spirí\textit{tui} \textbf{San}cto.

\item 12. Sicut erat in princípio, et nunc, et \textbf{sem}per,~* et in sǽcula sæcu\kern 0.02em\textit{lórum}. \textbf{A}men.
\end{itemize}

\smalltitle{II Sunday of Lent.}

\gscore[Ad Magnif.]{Ant. 1. f}{an_visionem_solesmes_1961}
\begin{otherlanguage}{english} \noindent Ant. Tell the vision that ye have seen to no man, until the Son of man be risen again from the dead.\end{otherlanguage}

\gresetinitiallines{0}
\gabcsnippet{(c4)1. Ma(f)gní(gh)fi(h)cat(h.) <v>\greheightstar</v>(:) á(h)ni(h)ma(h) <i>me</i>(g)<i>a</i>(f) <b>Dó</b>(gh)mi(gr)num.(gf..) (::) (Z-)
2. Et(f) ex(gh)sul(h)tá(h)vit(h) <b>spí</b>(ixi)ri(hr)tus(h) <b>me</b>(g hr)us(h.) <v>\greheightstar</v>(:) in(h) De(h)o(h) sa(h)lu(h)<i>tá</i>(g)<i>ri</i>(f) <b>me</b>(gh gr)o.(gf..) (::)}%%

\gresetinitiallines{1}

\begin{itemize}[label=\null, %
			leftmargin=0pt, %
			itemindent=0mm, %
			labelsep=0pt, %
			labelwidth=0pt, %
			rightmargin=0pt, %
			parsep=0pt, %
			topsep=0pt, %
			itemsep=0pt]%
%%Magnificat
%% each line requires \item

\item 3. Quia respéxit humilitátem an\textbf{cíl}læ \textbf{su}æ:~* ecce enim ex hoc beátam me dicent omnes gene\textit{rati}\textbf{ó}nes.

\item 4. Quia fecit mihi \textbf{ma}gna qui \textbf{pot}ens est:~* et sanctum \textit{nomen} \textbf{e}jus.

\item 5. Et misericórdia ejus a progénie \textbf{in} pro\textbf{gé}nies~* timén\textit{tibus} \textbf{e}um.

\item 6. Fecit poténtiam in \textbf{brá}chio \textbf{su}o:~* dispérsit supérbos mente \textit{cordis} \textbf{su}i.

\item 7. Depósuit po\textbf{tén}tes de \textbf{se}de,~* et exal\textit{távit} \textbf{hú}miles.

\item 8. Esuriéntes im\textbf{plé}vit \textbf{bo}nis:~* et dívites dimí\kern 0.02em\textit{sit} in\textbf{á}nes.

\item 9. Suscépit Israel \textbf{pú}erum \textbf{su}um,~* recordátus misericór\textit{diæ} \textbf{su}æ.

\item 10. Sicut locútus est ad \textbf{pa}tres \textbf{no}stros,~* Abraham et sémini e\kern 0.02em\textit{jus} \textit{in} \textbf{sǽ}cula.

\item 11. Glória \textbf{Pa}tri, et \textbf{Fí}lio,~* et Spirí\tinyhspace\textit{tui} \textbf{San}cto.

\item 12. Sicut erat in princípio, et \textbf{nunc,} et \textbf{sém}per, ~* et in sǽcula sæcu\tinyhspace\textit{lórum.} \textbf{A}men.

\end{itemize}

\smalltitle{III Sunday of Lent.}

\gscore[]{Ant. 8. G}{an_extollens_solesmes_1961}
\begin{otherlanguage}{english} \noindent Ant. A certain woman of the company lifted up her voice and said: Blessed is the womb that bare thee, and the paps which Thou hast sucked. But Jesus said unto her: Yea, rather, blessed are they that hear the word of God, and keep it.\end{otherlanguage}

\gresetinitiallines{0}

\gabcsnippet{(c4) 1. Ma(g)gní(hg)fi(gj)cat(j.) *(:) á(j)ni(j)ma(j) <i>me</i>(i)<i>a</i>(j) <b>Dó</b>(h)mi(gr)num.(g.) (::)}

\gabcsnippet{(c4) 2. Et(g) ex(h)sul(j)tá(j)vit(j) spí(j)ri(j)tus(j) <b>me</b>(k//jr)us(j.) *(:) in(j) De(j)o(j) sa(j)lu(j)<i>tá</i>(i)<i>ri</i>(j) <b>me</b>(h//gr)o.(g.) (::)}

\gresetinitiallines{1}
   \smallskip%
    \noindent%
    \begin{itemize}[%
    		label=\null, %
			leftmargin=0pt, %
			itemindent=0mm, %
			labelsep=0pt, %
			labelwidth=0pt, %
			rightmargin=0pt, %
			parsep=0pt, %
			topsep=0pt, %
			itemsep=0pt]%

\item 3. Quia respéxit humilitátem ancíllæ \textbf{su}æ:~* ecce enim ex hoc beátam me dicent omnes gene\textit{rati}\textbf{ó}nes.

\item 4. Quia fecit mihi magna qui \textbf{pot}ens est:~* et sanctum \textit{nomen} \textbf{e}jus.

\item 5. Et misericórdia ejus a progénie in pro\textbf{gé}nies~* timén\textit{tibus} \textbf{e}um.

\item 6. Fecit poténtiam in bráchio \textbf{su}o:~* dispérsit supérbos mente \textit{cordis} \textbf{su}i.

\item 7. Depósuit poténtes de \textbf{se}de,~* et exal\textit{távit} \textbf{hú}miles.

\item 8. Esuriéntes implévit \textbf{bo}nis:~* et dívites dimí\kern 0.04em\textit{sit} \textit{in}\textbf{á}nes.

\item 9. Suscépit Israel púerum \textbf{su}um,~* recordátus misericór\textit{diæ} \textbf{su}æ.

\item 10. Sicut locútus est ad patres \textbf{no}stros,~* Abraham et sémini e\kern 0.02em\textit{jus} \textit{in} \textbf{sǽ}cula.

\item 11. Glória Patri, et \textbf{Fí}lio,~* et Spirí\textit{tui} \textbf{San}cto.

\item 12. Sicut erat in princípio, et nunc, et \textbf{sem}per,~* et in sǽcula sæcu\kern 0.02em\textit{lórum}. \textbf{A}men.
\end{itemize}

\smalltitle{IV Sunday of Lent.}

\gscore[]{Ant. 1. g}{an_subiit_ergo_solesmes_1961}
\begin{otherlanguage}{english} \noindent Ant. And Jesus went up into a mountain, and there He sat with His disciples.\end{otherlanguage}

\gresetinitiallines{0}
\gabcsnippet{(c4) 1. Ma(f)gní(gh)fi(h)cat(h) * (:) á(h)ni(h)ma(h) <i>me</i>(g)a(f)
 <b>Dó</b>(gh)mi(gr)num.(g.) (::) (Z-)

2. Et(f) ex(gh)sul(h)tá(h)vit(h) <b>spí</b>(ixi)ri(hr)tus(h) <b>me</b>(g//hr)us(h.) *(:) in(h) De(h)o(h) sa(h)lu(h)<i>tá</i>(g)<i>ri</i>(f) <b>me</b>(gh//gr)o.(g.) (::)}

\gresetinitiallines{1}

\begin{itemize}[label=\null, %
			leftmargin=0pt, %
			itemindent=0mm, %
			labelsep=0pt, %
			labelwidth=0pt, %
			rightmargin=0pt, %
			parsep=0pt, %
			topsep=0pt, %
			itemsep=0pt]%
%%Magnificat
%% each line requires \item
\item 3. Quia respéxit humilitátem an\textbf{cíl}læ \textbf{su}æ:~* ecce enim ex hoc beátam me dicent omnes gene\textit{rati}\textbf{ó}nes.

\item 4. Quia fecit mihi \textbf{ma}gna qui \textbf{pot}ens est:~* et sanctum \textit{nomen} \textbf{e}jus.

\item 5. Et misericórdia ejus a progénie \textbf{in} pro\textbf{gé}nies~* timén\textit{tibus} \textbf{e}um.

\item 6. Fecit poténtiam in \textbf{brá}chio \textbf{su}o:~* dispérsit supérbos mente \textit{cordis} \textbf{su}i.

\item 7. Depósuit po\textbf{tén}tes de \textbf{se}de,~* et exal\textit{távit} \textbf{hú}miles.

\item 8. Esuriéntes im\textbf{plé}vit \textbf{bo}nis:~* et dívites dimí\kern 0.02em\textit{sit} in\textbf{á}nes.

\item 9. Suscépit Israel \textbf{pú}erum \textbf{su}um,~* recordátus misericór\textit{diæ} \textbf{su}æ.

\item 10. Sicut locútus est ad \textbf{pa}tres \textbf{no}stros,~* Abraham et sémini e\kern 0.02em\textit{jus} \textit{in} \textbf{sǽ}cula.

\item 11. Glória \textbf{Pa}tri, et \textbf{Fí}lio,~* et Spirí\tinyhspace\textit{tui} \textbf{San}cto.

\item 12. Sicut erat in princípio, et \textbf{nunc,} et \textbf{sém}per, ~* et in sǽcula sæcu\tinyhspace\textit{lórum.} \textbf{A}men.
\end{itemize}

\smalltitle{Passion Sunday.}

\gscore[]{Ant. 2. D}{an_abraham_pater_vester_solesmes_1961} 
\begin{otherlanguage}{english} \noindent Ant. Your father Abraham rejoiced to see My day and he saw it, and was glad.\end{otherlanguage}

\gresetinitiallines{0}

\gabcsnippet{(c3@f3)1. Ma(e)gní(fe)fi(eh)cat(h.) *(:) á(h)ni(h)ma(h) me(h)<i>a</i>(g) <b>Dó</b>(e)mi(fr)num.(f.) (::) (Z-)
2. Et<v>\hspace*{-0.3cm}\noindent</v>(e) ex(f)sul(h)tá(h)vit(h) spí(h)ri(h)tus(h) <b>me</b>(i//hr)us(h.) *(:) in(h) De(h)o(h) sa(h)lu(h)tá(h)<i>ri</i>(g) me(e//fr)o.(f.) (::)}

\gresetinitiallines{1}

\begin{itemize}[label=\null, %
			leftmargin=0pt, %
			itemindent=0mm, %
			labelsep=0pt, %
			labelwidth=0pt, %
			rightmargin=0pt, %
			parsep=0pt, %
			topsep=0pt, %
			itemsep=0pt]%
%%Magnificat
%% each line requires \item

\item 3. Quia respéxit humilitátem ancíllæ \textbf{su}æ:~* ecce enim ex hoc beátam me dicent omnes genera\kern 0.01 em\textit{ti}\textbf{ó}nes.

\item 4. Quia fecit mihi magna qui \textbf{pot}ens est:~* et sanctum no\textit{men} \textbf{e}jus.

\item 5. Et misericórdia ejus a progénie in pro\textbf{gé}nies~* timénti\textit{bus} \textbf{e}um.

\item 6. Fecit poténtiam in bráchio \textbf{su}o:~* dispérsit supérbos mente cor\textit{dis} \textbf{su}i.

\item 7. Depósuit poténtes de \textbf{se}de,~* et exaltá\textit{vit} \textbf{hú}miles.

\item 8. Esuriéntes implévit \textbf{bo}nis:~* et dívites dimísit \textit{in}\textbf{á}nes.

\item 9. Suscépit Israel púerum \textbf{su}um,~* recordátus misericórdi\kern 0.01 em\textit{æ} \textbf{su}æ.

\item 10. Sicut locútus est ad patres \textbf{no}stros,~* Abraham et sémini ejus \textit{in} \textbf{sǽ}cula.

\item 11. Glória Patri, et \textbf{Fí}lio,~* et Spirítu\kern 0.01 em\textit{i} \textbf{San}cto.

\item 12. Sicut erat in princípio, et nunc, et \textbf{sem}per,~* et in sǽcula sæculó\textit{rum}. \textbf{A}men.

\end{itemize}

\smalltitle{Palm Sunday.}

\gscore[]{Ant. 8. G*}{an_scriptum_est_enim__percutiam_solesmes_1961.1}

\begin{otherlanguage}{english} \noindent Ant. It is written: I will smite the Shepherd, and the sheep of the flock shall be scattered abroad, but after I am risen again, I will go before you into Galilee: there shall ye see Me, saith the Lord..\end{otherlanguage}

\rubrique{The tone is the same as on the I Sunday of Lent,\normaltext{ \pageref{1_sunday}.}}

\scripture{Luc. 1, 46–55.}

%\begin{paracol}{2}
%
%\lettrine{M}agníficat~* ánima mea Dóminum.\\ Et exsultávit spíritus meus~* in Deo salutári meo.
%
% Quia respéxit humilitátem ancíllæ suæ:~* ecce enim ex hoc beátam me dicent o\-mnes generatiónes.
%
% Quia fecit mihi magna qui potens est:~* et san\-ctum nomen ejus.
%
% Et misericórdia ejus a progénie in progénies~* timéntibus eum.
%
% Fecit poténtiam in bráchio suo:~* dispérsit supérbos mente cordis sui.
%
% Depósuit poténtes de sede,~* et exaltávit húmiles.
%
% Esuriéntes implévit bonis:~* et dívites dimísit inánes.
%
% Suscépit Israel púerum suum,~* recordátus misericórdiæ suæ.
%
% Sicut locútus est ad patres no\-stros,~* Abraham, et sémini ejus in sǽcula.
%
% Glória Patri, et Fílio,~* et Spirítui San\-cto.
%
% Sicut erat in princípio, et nunc, et semper,~* et in sǽcula sæculórum. Amen.
% 
% \switchcolumn

\begin{multicols}{2}
 \begin{otherlanguage}{english}
My soul \cc~* doth magnify the Lord.

And my spirit hath rejoiced~* in God my Saviour.

Because he hath regarded the humility of his handmaid;~* for behold from henceforth all generations shall call me blessed.

Because he that is mighty, hath done great things to me;~* and holy is his name.

And his mercy is from generation unto generations,~* to them that fear him.

He hath shewed might in his arm:~* he hath scattered the proud in the conceit of their heart.

He hath put down the mighty from their seat,~* and hath exalted the humble.

He hath filled the hungry with good things;~* and the rich he hath sent empty away.

He hath received Israel his servant,~* being mindful of his mercy:

As he spoke to our fathers,~* to Abraham and to his seed for ever.

Glory be to the Father, and to the Son,~* and to the Holy Ghost.

As it was in the beginning, is now,~* and ever shall be, world without end. Amen.
\end{otherlanguage}
% \end{paracol}
\end{multicols}

 \rubrique{\begin{otherlanguage}{english}
Sit as the antiphon is sung in full. Then the usual dialogue follows. On Sundays, all stand from now until the end of the office, unless one sits for the antiphons sung for the commemorations, according to custom.\end{otherlanguage}}

\textebilingue{dv}{dv_eng}
 \rubrique{\begin{otherlanguage}{english}
The celebrant sings the collect of the day, and any commemorations follow. When required, the Suffrage of All Saints is sung in the last place, as below.\end{otherlanguage}}

\bigtitle{Collect.}
\smallskip

\smalltitle{I Sunday of Lent.}

   \begin{paracol}{2}
   Orémus.
   
\lettrine{D}{e}us, qui Ecclésiam tuam ánnua quadragesimáli observatióne puríficas:~† præsta famíliæ tuæ; ut quod a te obtinére abstinéndo nítitur,* hoc bonis opéribus exsequátur. Per Dóminum.
   
   ℟. Amen.
   
\switchcolumn
\begin{otherlanguage}{english}
\noindent Let us pray.
O God, Who dost every year purge thy Church by the Fast of Forty Days, grant unto this thy family, that what things soever they strive to obtain at thy hand by abstaining from meats, they may ever turn to profit by good works.
Through Jesus Christ, thy Son our Lord, Who liveth and reigneth with thee, in the unity of the Holy Ghost, God, world without end.

 Amen.
\end{otherlanguage}
    \end{paracol}
    
    \smalltitle{II Sunday of Lent.}

   \begin{paracol}{2}
   Orémus.
   
\lettrine{D}{e}us, qui cónspicis omni nos virtúte destítui: intérius exteriúsque custódi;~† ut ab ómnibus adversitátibus muniámur in córpore,* et a právis cogitatiónibus mundémur in mente. Per Dóminum.
   
   ℟. Amen.
   
\switchcolumn
\begin{otherlanguage}{english}
\noindent Let us pray.
O God, Who seest that we have no power of ourselves to help ourselves, keep us both outwardly in our bodies, and inwardly in our souls, that we may be defended from all adversities which may happen to the body, and from all evil thoughts which may assault and hurt the soul.
Through Jesus Christ, thy Son our Lord, Who liveth and reigneth with thee, in the unity of the Holy Ghost, God, world without end.
 Amen.
\end{otherlanguage}
    \end{paracol}
    
    \smalltitle{III Sunday of Lent.}

   \begin{paracol}{2}
   Orémus.
   
 \lettrine[depth=1]{Q}{u}ǽsumus omnípotens Deus, vota humílium réspice:~† atque ad defensiónem nostram,* déxteram tuæ majestátis exténde.
Per Dóminum.
   
   ℟. Amen.
   
\switchcolumn
\begin{otherlanguage}{english}
\noindent Let us pray.
We beeseech thee, almighty God, look upon the hearty desires of thy humble servants, and stretch forth the right hand of thy Majesty to be our defense against all our enemies.
Through Jesus Christ, thy Son our Lord, Who liveth and reigneth with thee, in the unity of the Holy Ghost, God, world without end.
 Amen.
\end{otherlanguage}
    \end{paracol}
    
    \smalltitle{IV Sunday of Lent.}

   \begin{paracol}{2}
   Orémus.
   
\lettrine{C}{o}ncéde, quǽsumus omnípotens Deus:~† ut, qui ex mérito nostræ actiónis afflígimur,* tuæ grátiæ consolatióne respirémus.
Per Dóminum.
   
   ℟. Amen.
   
\switchcolumn
\begin{otherlanguage}{english}
\noindent Let us pray.
Grant, we beseech thee, Almighty God, that we who for our evil deeds are worthily punished, may, by the comfort of thy grace, mercifully be relieved.
Through Jesus Christ, thy Son our Lord, Who liveth and reigneth with thee, in the unity of the Holy Ghost, God, world without end.
 Amen.
\end{otherlanguage}
    \end{paracol}
    
    \smalltitle{Passion Sunday.}

   \begin{paracol}{2}
   Orémus.
   
\lettrine[depth=1]{Q}{u}ǽsumus omnípotens Deus, famíliam tuam propítius réspice:~† ut te largiénte, regátur in córpore;~* et te servánte, custodiátur in mente. Per Dóminum. 
   
   ℟. Amen.
   
\switchcolumn
\begin{otherlanguage}{english}
\noindent Let us pray.
We beseech thee, Almighty God, mercifully to look upon this thy family, that by thy great goodness they may be governed and preserved evermore, both in body and soul.
Through Jesus Christ, thy Son our Lord, Who liveth and reigneth with thee, in the unity of the Holy Ghost, God, world without end.
 Amen.
\end{otherlanguage}
    \end{paracol}
    
        \smalltitle{Palm Sunday.}

   \begin{paracol}{2}
   Orémus.
   
   \lettrine{O}{m}nípotens sempitérne Deus, qui humáno géneri, ad imitándum humilitátis exémplum, Salvatórem nostrum carnem súmere, et crucem subíre fecísti:~† concéde propítius; ut et patiéntiæ ipsíus habére documénta,*  et resurrectiónis consórtia mereámur. Per eúndem Dóminum nostrum.
   
   ℟. Amen.
   
\switchcolumn
\begin{otherlanguage}{english}
\noindent Let us pray.
Almighty and everlasting God, Who, of thy tender love towards mankind, hast sent thy Son our Saviour Jesus Christ to take upon Him our flesh and to suffer death upon the Cross, that all mankind should follow the example of His great humility; mercifully grant, that we may both follow the example of His patience, and also be made partakers of His resurrection.
Through the same Jesus Christ, thy Son, Our Lord, Who liveth and reigneth with thee in the unity of the Holy Ghost, God, world without end.
 Amen.
\end{otherlanguage}
    \end{paracol}

\myrule

\smalltitle{Suffrage of All Saints.}

\gscore[Ant.]{2.}{an_beata_dei_genitrix_virgo_solesmes_1961}
\begin{otherlanguage}{english} \noindent Ant. May the Blessed Virgin Mary, Mother of God, and all the Saints, intercede for us with the Lord. \end{otherlanguage}

\textebilingue{suffragium}{suffragium_en}

\textebilingue{dv}{dv_eng}

\rubrique{\begin{otherlanguage}{english}
The cantor sings the verse to which all reply.
\end{otherlanguage}}

\gscore[]{6.}{BD_Dominicis_Adventus_Quadragesimae}

\noindent \vv Let us bless the Lord.

\noindent \rr Thanks be to God.

\textebilingue{fidelium_animae}{fidelium_animae_en}

%\rubrique{\begin{otherlanguage}{english}
%Then one Our Father is said, unless Compline follows, or even some pious exercises, a sermon, or Benediction of the Blessed Sacrament, in parish and other smaller churches.\end{otherlanguage}}

%\rubrique{\begin{otherlanguage}{english}
%Then the office is concluded with the following verses, the Marian antiphon, etc.
%\end{otherlanguage}}


\end{document}